\documentclass[12pt]{article}

% ══════════════════════════════════════════════════════════════
%  CODIFICACIÓN E IDIOMA
% ══════════════════════════════════════════════════════════════
\usepackage[utf8]{inputenc}
\usepackage[T1]{fontenc}
\usepackage[spanish,es-nodecimaldot,es-tabla]{babel}
\usepackage[defaultfam,tabular,lining]{montserrat}
\spanishdecimal{.}

% ══════════════════════════════════════════════════════════════
%  FORMATO Y DISEÑO
% ══════════════════════════════════════════════════════════════
\usepackage[top=2.5cm, bottom=3.2cm, outer=3cm, inner=3cm]{geometry}
\usepackage{ragged2e}
\usepackage{setspace}
\usepackage{tocloft}
\usepackage{fancyhdr}
\usepackage{caption}
\usepackage{makeidx}
\makeindex
\usepackage{microtype}

% ══════════════════════════════════════════════════════════════
%  COLORES (PALETA GENERAL)
% ══════════════════════════════════════════════════════════════
\usepackage{xcolor}
\definecolor{blueunam}{RGB}{50, 70, 90}
\definecolor{catred}{RGB}{192,57,43}
\definecolor{anblue}{RGB}{41,128,185}
\definecolor{iupacgreen}{RGB}{39,100,55}
\definecolor{nerde}{RGB}{39,100,55}
\definecolor{separador}{RGB}{160,80,60}
\definecolor{atomo}{RGB}{128,0,128}
\definecolor{col1}{RGB}{210,228,255}
\definecolor{col2}{RGB}{200,238,210}
\definecolor{trans}{RGB}{255,238,200}
\definecolor{trans2}{RGB}{255,220,185}
\definecolor{pos}{RGB}{228,218,255}
\definecolor{especial}{RGB}{235,235,235}
\definecolor{rojoNota}{RGB}{180, 55, 55}
\definecolor{amarilloNota}{RGB}{200, 160, 50}
\definecolor{azulNota}{RGB}{55, 115, 165}

% ══════════════════════════════════════════════════════════════
%  MACROS DE COLOR (CORREGIDOS)
% ══════════════════════════════════════════════════════════════
% \ac: Átomo central en morado con tipografía recta (romana)
\newcommand{\ac}[1]{{\color{atomo}\textnormal{#1}}}
\newcommand{\carga}[1]{^{#1}}

% ══════════════════════════════════════════════════════════════
%  MATEMÁTICAS
% ══════════════════════════════════════════════════════════════
\usepackage{amsmath}
\usepackage{mathtools}
\usepackage{amsthm}
\usepackage{amssymb}
\usepackage{amsbsy}
\usepackage{amsfonts}
\usepackage{amstext}
\usepackage{amscd}
\usepackage{nccmath}
\usepackage{bbm}
\usepackage{esint}
\usepackage{cancel}
\usepackage{braket}
\usepackage{tensor}
\setcounter{MaxMatrixCols}{30}
\allowdisplaybreaks

% ══════════════════════════════════════════════════════════════
%  UNIDADES
% ══════════════════════════════════════════════════════════════
\usepackage{siunitx}
\sisetup{output-decimal-marker={.}, group-separator={,}}

% ══════════════════════════════════════════════════════════════
%  GRÁFICOS Y FIGURAS
% ══════════════════════════════════════════════════════════════
\usepackage{graphicx}
\usepackage{float}
\usepackage{tikz}
\AtBeginDocument{\shorthandoff{>}}
\usetikzlibrary{shapes.geometric, arrows.meta, positioning, decorations.markings}
\usepackage{pgfplots}
\pgfplotsset{compat=1.18}

% ══════════════════════════════════════════════════════════════
%  TABLAS
% ══════════════════════════════════════════════════════════════
\usepackage{array}
\usepackage{tabularx}
\usepackage{multirow}
\usepackage{booktabs}
\usepackage{longtable}
\usepackage{colortbl}
\usepackage{pdflscape}
\usepackage{multicol}

% ══════════════════════════════════════════════════════════════
%  QUÍMICA
% ══════════════════════════════════════════════════════════════
\usepackage{chemfig}
\usepackage[version=4]{mhchem}
\usepackage{elements}
\usepackage{chemnum}
\usepackage{chemplants}
\usepackage{carbohydrates}
\usepackage{chemobabel}
\usepackage{tikzorbital}
\usepackage{modiagram}

% ══════════════════════════════════════════════════════════════
%  DIAMANTE NFPA 704
% ══════════════════════════════════════════════════════════════
\newcommand{\NFPAdiamond}[4]{%
    \begin{tikzpicture}[scale=0.45, every node/.style={transform shape}]
        \draw[thick, line width=1.5pt] (0,2) -- (2,4) -- (4,2) -- (2,0) -- cycle;
        \draw[thick] (0,2) -- (4,2);
        \draw[thick] (2,0) -- (2,4);
        \fill[red]    (2,4) -- (1,3) -- (2,2) -- (3,3) -- cycle;
        \node[text=white, font=\bfseries\Large] at (2,3.1)  {#2};
        \fill[blue]   (0,2) -- (1,3) -- (2,2) -- (1,1) -- cycle;
        \node[text=white, font=\bfseries\Large] at (0.9,2)  {#1};
        \fill[yellow] (4,2) -- (3,3) -- (2,2) -- (3,1) -- cycle;
        \node[text=black, font=\bfseries\Large] at (3.1,2)  {#3};
        \fill[white]  (2,0) -- (1,1) -- (2,2) -- (3,1) -- cycle;
        \node[text=black, font=\bfseries\large] at (2,0.9)  {#4};
        \draw[thick, line width=1.5pt] (0,2) -- (2,4) -- (4,2) -- (2,0) -- cycle;
    \end{tikzpicture}%
}

% ══════════════════════════════════════════════════════════════
%  LISTAS
% ══════════════════════════════════════════════════════════════
\usepackage[shortlabels]{enumitem}

% ══════════════════════════════════════════════════════════════
%  CAJAS Y RESALTADO
% ══════════════════════════════════════════════════════════════
\usepackage[most]{tcolorbox}

% ══════════════════════════════════════════════════════════════
%  NOTAS POR NIVEL
% ══════════════════════════════════════════════════════════════
\newtcolorbox{notafija}{
    enhanced, breakable, colback=white, colframe=rojoNota, coltext=black,
    boxrule=0pt, leftrule=3pt, arc=0pt, left=8pt, right=8pt, top=6pt, bottom=6pt, fontupper=\small
}
\newtcolorbox{notaimportante}{
    enhanced, breakable, colback=white, colframe=amarilloNota, coltext=black,
    boxrule=0pt, leftrule=2pt, arc=0pt, left=8pt, right=8pt, top=6pt, bottom=6pt, fontupper=\small\itshape
}
\newtcolorbox{notadidactica}{
    enhanced, breakable, colback=white, colframe=azulNota, coltext=black!70,
    boxrule=0pt, leftrule=1pt, arc=0pt, left=8pt, right=8pt, top=4pt, bottom=4pt, fontupper=\small\itshape
}

% ══════════════════════════════════════════════════════════════
%  ENCABEZADO Y PIE
% ══════════════════════════════════════════════════════════════
\pagestyle{fancy}
\fancyhf{}
\fancyfoot[R]{\thepage}
\renewcommand{\headrulewidth}{0pt}

\newcommand{\activarEncabezado}[1]{%
  \fancyhead[L]{\small\textcolor{blueunam}{#1}}%
  \fancyhead[R]{\small\textcolor{blueunam}{\nouppercase{\leftmark}}}%
  \renewcommand{\headrulewidth}{0.4pt}%
  \renewcommand{\headrule}{\hbox to\headwidth{\color{blueunam}\leaders\hrule height \headrulewidth\hfill}}%
}

\renewcommand*\oldstylenums[1]{{\fontfamily{Montserrat-TOsF}\selectfont #1}}

% ══════════════════════════════════════════════════════════════
%  ÍNDICE Y BIBLIOGRAFÍA
% ══════════════════════════════════════════════════════════════
\AtBeginDocument{
    \renewcommand{\contentsname}{\centering\fontsize{12}{18}\selectfont\textcolor{blueunam}{\textbf{ÍNDICE}}}
}
\AtBeginDocument{
    \renewcommand{\bibname}{\centering\fontsize{12}{18}\selectfont\textcolor{blueunam}{\textbf{REFERENCIAS}}}
}
\renewcommand{\cftsecfont}{\color{blueunam}}
\setcounter{tocdepth}{3}
\setcounter{secnumdepth}{3}

% ══════════════════════════════════════════════════════════════
%  COLUMNAS PERSONALIZADAS
% ══════════════════════════════════════════════════════════════
\newcolumntype{L}[1]{>{\raggedright\arraybackslash}m{#1}}
\newcolumntype{C}[1]{>{\centering\arraybackslash}m{#1}}

% ══════════════════════════════════════════════════════════════
%  CELDAS DE TABLA DE CATIONES (Compacta)
% ══════════════════════════════════════════════════════════════
\newcommand{\catcell}[4]{%
    \cellcolor{#1}%
    \begin{minipage}[t][1.4cm][t]{1.3cm}%
        \centering\vspace{2pt}%
        {\scriptsize\textbf{#2}}\par\vspace{1pt}%
        {\tiny\color{gray}#3}\par\vspace{1pt}%
        {\tiny\color{blueunam}#4}%
    \end{minipage}}

\newcommand{\emptycell}{%
    \cellcolor{white}%
    \begin{minipage}[t][1.4cm][t]{1.3cm}\end{minipage}}

% ══════════════════════════════════════════════════════════════
%  AMBIENTES DE EJERCICIO
% ══════════════════════════════════════════════════════════════
\newtcbtheorem{exercise}{Ejercicio}{
    enhanced, breakable,
    colback=gray!10, colframe=gray!50,
    boxrule=0.4pt, arc=0pt, left=0.2cm,
    attach boxed title to top left={yshift=-2mm,xshift=-2mm},
    coltitle=black, fonttitle=\bfseries,
    colbacktitle=white,
    boxed title style={boxrule=.4pt,sharp corners, colframe=gray!50}
}{exercise}

\newtcbtheorem{exerciseline}{Ejercicio}{
    enhanced, frame empty, interior empty,
    colframe=blueunam,
    borderline west={1pt}{0pt}{blueunam},
    left=0.2cm,
    attach boxed title to top left={yshift=-2mm,xshift=-2mm},
    coltitle=black, fonttitle=\bfseries,
    colbacktitle=white,
    boxed title style={boxrule=.4pt,sharp corners}
}{exerciseline}

% ══════════════════════════════════════════════════════════════
%  REFERENCIAS Y CITAS
% ══════════════════════════════════════════════════════════════
\usepackage{csquotes}
\usepackage{url}
\usepackage{xurl}
\usepackage{fancyvrb}
\usepackage{witharrows}
\usepackage[hidelinks]{hyperref}
\usepackage{cleveref}
\usepackage[backend=biber, style=apa, language=spanish]{biblatex}

% ══════════════════════════════════════════════════════════════
%  PIES DE FIGURA Y TABLA (APA)
% ══════════════════════════════════════════════════════════════
\DeclareCaptionLabelSeparator*{spaced}{\\[2ex]}
\captionsetup[table]{textfont=it,format=plain,justification=justified,
  singlelinecheck=false,labelsep=spaced,skip=0pt}
\captionsetup[figure]{labelsep=period,labelfont=bf,justification=justified,
  singlelinecheck=false,font=doublespacing}

\setlength{\parindent}{0pt}

% \addbibresource{referencias_NomenclaturaIno.bib}

% ══════════════════════════════════════════════════════════════
%  DOCUMENTO
% ══════════════════════════════════════════════════════════════
\begin{document}

%---------- Portada ----------%
\begin{titlepage}
\thispagestyle{empty}
\begin{center}
\vspace*{2cm}
\fontsize{14}{30}\selectfont\textbf{Universidad Nacional Autónoma de México} \\
\vspace{1cm}
\fontsize{14}{30}\selectfont\textbf{Química e Ingeniería en Materiales} \\
\vspace{1cm}
\fontsize{14}{30}\selectfont\textbf{Nomenclatura Inorgánica} \\
\end{center}
\vspace{2cm}
\begin{flushright}
    \fontsize{14}{30}\selectfont
        \hspace{2cm} Nombre: Uriel Sánchez Méndez \\
        \hspace{2cm} Matrícula: 322647576 \\
        \hspace{2cm} Apuntes de Nomenclatura Inorgánica \\
        \hspace{2cm} Fecha: \today \\
\end{flushright}
\end{titlepage}

%---------- Tabla de contenido ----------%
\tableofcontents
\thispagestyle{empty}
\newpage

%---------- Contenido ----------%
\spacing{1.5}
% Configuración global de tablas
\setlength{\tabcolsep}{4pt}
\renewcommand{\arraystretch}{1.3}

%---------- NOTA INTRODUCTORIA ----------%
\phantomsection
\addcontentsline{toc}{section}{Nota sobre los sistemas de nomenclatura}
\section*{\centering\fontsize{12}{18}\selectfont{\textcolor{blueunam}{\textbf{NOTA SOBRE LOS SISTEMAS DE NOMENCLATURA}}}}

En química inorgánica existen formalmente tres sistemas de nomenclatura:

\begin{itemize}
  \item \textbf{Stock:} indica el estado de oxidación del catión mediante número romano entre paréntesis. No modifica el nombre del anión. Ejemplo: \ce{FeCl2} = cloruro de hierro(II).
  \item \textbf{Tradicional (clásica):} usa sufijos y prefijos tanto para el catión (-oso/-ico) como para el anión (-ito/-ato, hipo-, per-). Ejemplo: \ce{FeCl2} = cloruro ferroso.
  \item \textbf{IUPAC:} indica la carga de cada ion entre paréntesis con números arábigos y signo después. Ejemplo: \ce{FeCl2} = cloruro(1$-$) de hierro(2+).
\end{itemize}

\textbf{El problema real:} ningún profesor, libro de texto o laboratorio utiliza ninguno de estos tres sistemas de forma pura. Lo que se usa en la práctica es un sistema híbrido: el anión se nombra en tradicional (hipoclorito, sulfato, nitrito) y el catión se nombra con número romano cuando tiene más de un estado de oxidación (hierro(II), cobre(II)), pero se omite el romano cuando solo tiene uno (sodio, aluminio, plata).

Este sistema híbrido no tiene nombre oficial en ningún libro. Se le llama aquí \textbf{Usual} porque es el nombre más honesto para lo que realmente se practica, sin pretender que es Stock puro ni Tradicional puro. En la mayoría de libros de texto y en la práctica universitaria (Chang, Petrucci, apuntes UNAM) este sistema híbrido se conoce simplemente como \textbf{Stock}, aunque técnicamente Stock puro usaría el romano también para los aniones. La confusión persiste porque los libros enseñan los sistemas formales pero los aplican mezclados, sin admitirlo.

Las tablas de este documento usan tres columnas: \textbf{Usual} (el sistema híbrido real), \textbf{IUPAC} (el sistema formal internacional) e \textbf{Inglés} (nombre en inglés según IUPAC).

Para las reglas completas y oficiales de nomenclatura inorgánica, ver:\\
\textbf{Libro Rojo IUPAC:} Nomenclature of Inorganic Chemistry. IUPAC Recommendations 2005.\\
Disponible en: \url{https://iupac.org/what-we-do/books/redbook/}

\newpage

%---------- REGLAS DE CONSTRUCCIÓN ----------%
\phantomsection
\addcontentsline{toc}{section}{Reglas de Construcción de Nombres}
\section*{\centering\fontsize{12}{18}\selectfont{\textcolor{blueunam}{\textbf{REGLAS DE CONSTRUCCIÓN DE NOMBRES}}}}

\subsection*{\textcolor{blueunam}{Sistema Usual}}

\textbf{Aniones monoatómicos:} raíz del elemento + \textbf{-uro}. Excepción: \ce{O^{2-}} = óxido (no oxuro).

\textbf{Oxoaniones:} el nombre depende del número de oxígenos respecto al anión de referencia (-ato):
\begin{itemize}
  \item Un O menos que -ato: \textbf{hipo-\ldots-ito}
  \item -ato menos un O: \textbf{-ito}
  \item Referencia: \textbf{-ato}
  \item Un O más que -ato: \textbf{per-\ldots-ato}
\end{itemize}
Cada paso sube el E.O. del átomo central en 2 unidades. El patrón completo solo aplica a halógenos; azufre, nitrógeno y fósforo solo tienen -ito y -ato.

\textbf{Cationes:} nombre del elemento. Se agrega número romano entre paréntesis \textbf{solo si el elemento tiene más de un estado de oxidación común}. Si tiene uno solo, se omite el romano: \ce{Na+} = sodio, \ce{Al^{3+}} = aluminio, \ce{Fe^{2+}} = hierro(II).

\textbf{Sales y compuestos:} anión primero + ``de'' + catión. Ejemplo: \ce{Fe2(SO4)3} = sulfato de hierro(III).

\textbf{Oxácidos:} -ato del anión $\to$ ácido \ldots-ico; \quad -ito del anión $\to$ ácido \ldots-oso.\\
Ejemplo: nitrato $\to$ ácido nítrico; nitrito $\to$ ácido nitroso.

\textbf{Hidrácidos:} en disolución acuosa: ácido + raíz + \textbf{-hídrico}.\\
Ejemplo: \ce{HCl}(aq) = ácido clorhídrico.

\vspace{0.5em}

\subsection*{\textcolor{blueunam}{Sistema IUPAC}}

\textbf{Prefijos para número de átomos:}

\renewcommand{\arraystretch}{1.2}
\begin{longtable}{C{1.4cm} C{1.8cm} C{1.4cm} C{1.8cm} C{1.4cm} C{1.8cm} C{1.4cm} C{1.8cm}}
\toprule
\textbf{N} & \textbf{Prefijo} & \textbf{N} & \textbf{Prefijo} & \textbf{N} & \textbf{Prefijo} & \textbf{N} & \textbf{Prefijo} \\
\midrule
1 & mono* & 3 & tri  & 5 & penta & 7 & hepta \\
2 & di    & 4 & tetra& 6 & hexa  & 8 & octa  \\
\bottomrule
\end{longtable}
{\small *mono se omite cuando hay un solo átomo del elemento central. No se omite para los oxígenos si hay uno solo: \ce{ClO-} = oxoclorato(1$-$), no monooxoclorato.}

\vspace{0.4em}

\textbf{Aniones monoatómicos:} nombre del elemento + carga entre paréntesis con signo después.\\
\ce{Cl-} = cloruro(1$-$); \quad \ce{O^{2-}} = óxido(2$-$); \quad \ce{Fe^{3+}} = hierro(3+).

\textbf{Oxoaniones:} prefijo oxo $\times$ número de O + raíz del elemento central + \textbf{-ato} + carga entre paréntesis.\\
Construcción paso a paso:
\begin{enumerate}
  \item Contar los oxígenos: $n$ oxígenos $\to$ prefijo griego + ``oxo''
  \item Raíz del elemento central
  \item Terminación \textbf{-ato} siempre
  \item Carga del ion entre paréntesis: número arábigo + signo después
\end{enumerate}
Ejemplo: \ce{SO4^{2-}} $\to$ 4 O = tetraoxo + sulfato + (2$-$) = \textbf{tetraoxosulfato(2$-$)}.\\
Ejemplo con dos metales: \ce{Cr2O7^{2-}} $\to$ 7 O + 2 Cr = heptaoxo\textbf{di}cromato(2$-$).

\textbf{Oxácidos IUPAC:} ácido + prefijos oxo + raíz + \textbf{-ico} + E.O. del átomo central en romano.\\
Ejemplo: \ce{H2SO4}: S en 6+ $\to$ 4 O $\to$ \textbf{ácido tetraoxosulfúrico(VI)}.

\textbf{Diferencia clave:} en oxoaniones la carga es la del \textit{ion}; en oxácidos el paréntesis indica el E.O. del átomo central en romano.

\textbf{Sales IUPAC:} anión(carga) + ``de'' + catión(carga).\\
Ejemplo: \ce{FeCl2} = cloruro(1$-$) de hierro(2+).

\newpage

%---------- OXOANIONES ----------%
\phantomsection
\addcontentsline{toc}{section}{Oxoaniones}
\section*{\centering\fontsize{12}{18}\selectfont{\textcolor{blueunam}{\textbf{OXOANIONES}}}}

Los oxoaniones son aniones poliatómicos que contienen oxígeno. El patrón para halógenos es:
\textbf{hipo$\cdots$ito} $\to$ \textbf{$\cdots$ito} $\to$ \textbf{$\cdots$ato} $\to$ \textbf{per$\cdots$ato},
donde cada paso agrega un O y sube el E.O. en 2 unidades.

\renewcommand{\arraystretch}{1.3}
\setlength{\tabcolsep}{5pt}

\vspace{0.3em}
\textbf{Halógenos ({\color{atomo}Cl, Br, I})}

\newgeometry{top=2cm, bottom=2cm, left=1.5cm, right=1.5cm}
\begin{longtable}{C{2.2cm} >{\color{catred}}C{1.4cm} L{4.0cm} >{\color{iupacgreen}}L{6.2cm} L{2.8cm}}
\rule{0pt}{2.8ex}\textbf{Fórmula} & \textcolor{nerde}{\textbf{E.O.}} & \textcolor{nerde}{\textbf{Usual}} & \textcolor{nerde}{\textbf{IUPAC}} & \textcolor{nerde}{\textbf{Inglés}} \\
\midrule
\endfirsthead
\rule{0pt}{2.8ex}\textbf{Fórmula} & \textcolor{nerde}{\textbf{E.O.}} & \textcolor{nerde}{\textbf{Usual}} & \textcolor{nerde}{\textbf{IUPAC}} & \textcolor{nerde}{\textbf{Inglés}} \\
\midrule
\endhead
\endlastfoot
\ce{\ac{Cl}O^-} & 1+ & Hipoclorito & Oxoclorato(1$-$)       & Hypochlorite \\
\ce{\ac{Cl}O2^-} & 3+ & Clorito     & Dioxoclorato(1$-$)     & Chlorite \\
\ce{\ac{Cl}O3^-} & 5+ & Clorato     & Trioxoclorato(1$-$)    & Chlorate \\
\ce{\ac{Cl}O4^-} & 7+ & Perclorato  & Tetraoxoclorato(1$-$)  & Perchlorate \\
\arrayrulecolor{separador}\midrule\arrayrulecolor{black}
\ce{\ac{Br}O^-} & 1+ & Hipobromito & Oxobromato(1$-$)       & Hypobromite \\
\ce{\ac{Br}O2^-} & 3+ & Bromito     & Dioxobromato(1$-$)     & Bromite \\
\ce{\ac{Br}O3^-} & 5+ & Bromato     & Trioxobromato(1$-$)    & Bromate \\
\ce{\ac{Br}O4^-} & 7+ & Perbromato  & Tetraoxobromato(1$-$)  & Perbromate \\
\arrayrulecolor{separador}\midrule\arrayrulecolor{black}
\ce{\ac{I}O^-} & 1+ & Hipoyodito  & Oxoyodato(1$-$)        & Hypoiodite \\
\ce{\ac{I}O2^-} & 3+ & Yodito      & Dioxoyodato(1$-$)      & Iodite \\
\ce{\ac{I}O3^-} & 5+ & Yodato      & Trioxoyodato(1$-$)     & Iodate \\
\ce{\ac{I}O4^-} & 7+ & Peryodato   & Tetraoxoyodato(1$-$)   & Periodate \\
\end{longtable}

\vspace{0.3em}
\textbf{Anfígenos ({\color{atomo}S, Se, Te})}

\begin{longtable}{C{2.2cm} >{\color{catred}}C{1.4cm} L{4.0cm} >{\color{iupacgreen}}L{6.2cm} L{2.8cm}}
\rule{0pt}{2.8ex}\textbf{Fórmula} & \textcolor{nerde}{\textbf{E.O.}} & \textcolor{nerde}{\textbf{Usual}} & \textcolor{nerde}{\textbf{IUPAC}} & \textcolor{nerde}{\textbf{Inglés}} \\
\midrule\endhead\endlastfoot
\ce{\ac{S}O3^{2-}} & 4+ & Sulfito   & Trioxosulfato(2$-$)    & Sulfite \\
\ce{\ac{S}O4^{2-}} & 6+ & Sulfato   & Tetraoxosulfato(2$-$)  & Sulfate \\
\arrayrulecolor{separador}\midrule\arrayrulecolor{black}
\ce{\ac{Se}O3^{2-}} & 4+ & Selenito  & Trioxoselenato(2$-$)   & Selenite \\
\ce{\ac{Se}O4^{2-}} & 6+ & Selenato  & Tetraoxoselenato(2$-$) & Selenate \\
\arrayrulecolor{separador}\midrule\arrayrulecolor{black}
\ce{\ac{Te}O3^{2-}} & 4+ & Telurito  & Trioxotelurato(2$-$)   & Tellurite \\
\ce{\ac{Te}O4^{2-}} & 6+ & Telurato  & Tetraoxotelurato(2$-$) & Tellurate \\
\end{longtable}

\vspace{0.3em}
\textbf{Nitrogenoides ({\color{atomo}N, P, As, Sb})}

\begin{longtable}{C{2.2cm} >{\color{catred}}C{1.4cm} L{4.0cm} >{\color{iupacgreen}}L{6.2cm} L{2.8cm}}
\rule{0pt}{2.8ex}\textbf{Fórmula} & \textcolor{nerde}{\textbf{E.O.}} & \textcolor{nerde}{\textbf{Usual}} & \textcolor{nerde}{\textbf{IUPAC}} & \textcolor{nerde}{\textbf{Inglés}} \\
\midrule\endhead\endlastfoot
\ce{\ac{N}O2^-} & 3+ & Nitrito      & Dioxonitrato(1$-$)        & Nitrite \\
\ce{\ac{N}O3^-} & 5+ & Nitrato      & Trioxonitrato(1$-$)       & Nitrate \\
\arrayrulecolor{separador}\midrule\arrayrulecolor{black}
\ce{\ac{P}O3^{3-}} & 3+ & Fosfito      & Trioxofosfato(3$-$)       & Phosphite \\
\ce{\ac{P}O4^{3-}} & 5+ & Fosfato      & Tetraoxofosfato(3$-$)     & Phosphate \\
\arrayrulecolor{separador}\midrule\arrayrulecolor{black}
\ce{\ac{As}O3^{3-}} & 3+ & Arsenito     & Trioxoarseniato(3$-$)     & Arsenite \\
\ce{\ac{As}O4^{3-}} & 5+ & Arseniato    & Tetraoxoarseniato(3$-$)   & Arsenate \\
\arrayrulecolor{separador}\midrule\arrayrulecolor{black}
\ce{\ac{Sb}O3^{3-}} & 3+ & Antimonito   & Trioxoantimoniato(3$-$)   & Antimonite \\
\ce{\ac{Sb}O4^{3-}} & 5+ & Antimoniato  & Tetraoxoantimoniato(3$-$) & Antimonate \\
\end{longtable}

\vspace{0.3em}
\textbf{Carbonoideos ({\color{atomo}C, Si})}

\begin{longtable}{C{2.2cm} >{\color{catred}}C{1.4cm} L{4.0cm} >{\color{iupacgreen}}L{6.2cm} L{2.8cm}}
\rule{0pt}{2.8ex}\textbf{Fórmula} & \textcolor{nerde}{\textbf{E.O.}} & \textcolor{nerde}{\textbf{Usual}} & \textcolor{nerde}{\textbf{IUPAC}} & \textcolor{nerde}{\textbf{Inglés}} \\
\midrule\endhead\endlastfoot
\ce{\ac{C}O3^{2-}} & 4+ & Carbonato    & Trioxocarbonato(2$-$)     & Carbonate \\
\arrayrulecolor{separador}\midrule\arrayrulecolor{black}
\ce{\ac{Si}O3^{2-}} & 4+ & Silicato     & Trioxosilicato(2$-$)      & Silicate \\
\ce{\ac{Si}O4^{4-}} & 4+ & Ortosilicato & Tetraoxosilicato(4$-$)    & Orthosilicate \\
\end{longtable}

\vspace{0.3em}
\textbf{Metales de transición ({\color{atomo}Cr, Mn, V, Mo, W})}

\begin{longtable}{C{2.2cm} >{\color{catred}}C{1.4cm} L{4.0cm} >{\color{iupacgreen}}L{6.2cm} L{2.8cm}}
\rule{0pt}{2.8ex}\textbf{Fórmula} & \textcolor{nerde}{\textbf{E.O.}} & \textcolor{nerde}{\textbf{Usual}} & \textcolor{nerde}{\textbf{IUPAC}} & \textcolor{nerde}{\textbf{Inglés}} \\
\midrule\endhead\endlastfoot
\ce{\ac{V}O3^-} & 5+ & Vanadato      & Trioxovanadato(1$-$)       & Vanadate \\
\arrayrulecolor{separador}\midrule\arrayrulecolor{black}
\ce{\ac{Cr}O4^{2-}} & 6+ & Cromato       & Tetraoxocromato(2$-$)      & Chromate \\
\ce{\ac{Cr}2O7^{2-}} & 6+ & Dicromato     & Heptaoxodicromato(2$-$)    & Dichromate \\
\arrayrulecolor{separador}\midrule\arrayrulecolor{black}
\ce{\ac{Mn}O4^-} & 7+ & Permanganato  & Tetraoxomanganato(1$-$)    & Permanganate \\
\ce{\ac{Mn}O4^{2-}} & 6+ & Manganato     & Tetraoxomanganato(2$-$)    & Manganate \\
\arrayrulecolor{separador}\midrule\arrayrulecolor{black}
\ce{\ac{Mo}O4^{2-}} & 6+ & Molibdato     & Tetraoxomolibdato(2$-$)    & Molybdate \\
\arrayrulecolor{separador}\midrule\arrayrulecolor{black}
\ce{\ac{W}O4^{2-}} & 6+ & Tungstato     & Tetraoxotungstato(2$-$)    & Tungstate \\
\end{longtable}

\newpage

%---------- ANIONES MONOATÓMICOS ----------%
\phantomsection
\addcontentsline{toc}{section}{Aniones Monoatómicos}
\section*{\centering\fontsize{12}{18}\selectfont{\textcolor{blueunam}{\textbf{ANIONES MONOATÓMICOS}}}}

Terminación característica: \textbf{-uro}. Excepción: oxígeno forma \textbf{óxido} (no oxuro).

\vspace{0.3em}
\textbf{Halógenos ({\color{atomo}F, Cl, Br, I})}

\begin{longtable}{C{2.2cm} >{\color{catred}}C{1.4cm} L{4.0cm} >{\color{iupacgreen}}L{6.2cm} L{2.8cm}}
\rule{0pt}{2.8ex}\textbf{Fórmula} & \textcolor{nerde}{\textbf{E.O.}} & \textcolor{nerde}{\textbf{Usual}} & \textcolor{nerde}{\textbf{IUPAC}} & \textcolor{nerde}{\textbf{Inglés}} \\
\midrule\endhead\endlastfoot
\ce{\ac{F}^-} & 1$-$ & Fluoruro & Fluoruro(1$-$) & Fluoride \\
\arrayrulecolor{separador}\midrule\arrayrulecolor{black}
\ce{\ac{Cl}^-} & 1$-$ & Cloruro  & Cloruro(1$-$)  & Chloride \\
\arrayrulecolor{separador}\midrule\arrayrulecolor{black}
\ce{\ac{Br}^-} & 1$-$ & Bromuro  & Bromuro(1$-$)  & Bromide \\
\arrayrulecolor{separador}\midrule\arrayrulecolor{black}
\ce{\ac{I}^-} & 1$-$ & Yoduro   & Yoduro(1$-$)   & Iodide \\
\end{longtable}

\vspace{0.3em}
\textbf{Anfígenos ({\color{atomo}O, S, Se, Te}) + Peróxido (\ce{\ac{O}2^{2-}})}

\begin{longtable}{C{2.2cm} >{\color{catred}}C{1.4cm} L{4.0cm} >{\color{iupacgreen}}L{6.2cm} L{2.8cm}}
\rule{0pt}{2.8ex}\textbf{Fórmula} & \textcolor{nerde}{\textbf{E.O.}} & \textcolor{nerde}{\textbf{Usual}} & \textcolor{nerde}{\textbf{IUPAC}} & \textcolor{nerde}{\textbf{Inglés}} \\
\midrule\endhead\endlastfoot
\ce{\ac{O}^{2-}} & 2$-$     & Óxido    & Óxido(2$-$)    & Oxide \\
\ce{\ac{O}2^{2-}} & 1$-$/cad & Peróxido & Peróxido(2$-$) & Peroxide \\
\arrayrulecolor{separador}\midrule\arrayrulecolor{black}
\ce{\ac{S}^{2-}} & 2$-$     & Sulfuro    & Sulfuro(2$-$)    & Sulfide \\
\arrayrulecolor{separador}\midrule\arrayrulecolor{black}
\ce{\ac{Se}^{2-}} & 2$-$     & Selenuro   & Selenuro(2$-$)   & Selenide \\
\arrayrulecolor{separador}\midrule\arrayrulecolor{black}
\ce{\ac{Te}^{2-}} & 2$-$     & Telururo   & Telururo(2$-$)   & Telluride \\
\end{longtable}

\vspace{0.3em}
\textbf{Nitrogenoides ({\color{atomo}N, P, As, Sb})}

\begin{longtable}{C{2.2cm} >{\color{catred}}C{1.4cm} L{4.0cm} >{\color{iupacgreen}}L{6.2cm} L{2.8cm}}
\rule{0pt}{2.8ex}\textbf{Fórmula} & \textcolor{nerde}{\textbf{E.O.}} & \textcolor{nerde}{\textbf{Usual}} & \textcolor{nerde}{\textbf{IUPAC}} & \textcolor{nerde}{\textbf{Inglés}} \\
\midrule\endhead\endlastfoot
\ce{\ac{N}^{3-}} & 3$-$ & Nitruro     & Nitruro(3$-$)     & Nitride \\
\arrayrulecolor{separador}\midrule\arrayrulecolor{black}
\ce{\ac{P}^{3-}} & 3$-$ & Fosfuro     & Fosfuro(3$-$)     & Phosphide \\
\arrayrulecolor{separador}\midrule\arrayrulecolor{black}
\ce{\ac{As}^{3-}} & 3$-$ & Arseniuro   & Arseniuro(3$-$)   & Arsenide \\
\arrayrulecolor{separador}\midrule\arrayrulecolor{black}
\ce{\ac{Sb}^{3-}} & 3$-$ & Antimoniuro & Antimoniuro(3$-$) & Antimonide \\
\end{longtable}

\vspace{0.3em}
\textbf{Carbonoideos ({\color{atomo}C, Si})}

\begin{longtable}{C{2.2cm} >{\color{catred}}C{1.4cm} L{4.0cm} >{\color{iupacgreen}}L{6.2cm} L{2.8cm}}
\rule{0pt}{2.8ex}\textbf{Fórmula} & \textcolor{nerde}{\textbf{E.O.}} & \textcolor{nerde}{\textbf{Usual}} & \textcolor{nerde}{\textbf{IUPAC}} & \textcolor{nerde}{\textbf{Inglés}} \\
\midrule\endhead\endlastfoot
\ce{\ac{C}^{4-}} & 4$-$ & Carburo   & Carburo(4$-$)   & Carbide \\
\arrayrulecolor{separador}\midrule\arrayrulecolor{black}
\ce{\ac{Si}^{4-}} & 4$-$ & Siliciuro & Siliciuro(4$-$) & Silicide \\
\end{longtable}

\vspace{0.3em}
\textbf{Hidrógeno (H)}

\begin{longtable}{C{2.2cm} >{\color{catred}}C{1.4cm} L{4.0cm} >{\color{iupacgreen}}L{6.2cm} L{2.8cm}}
\rule{0pt}{2.8ex}\textbf{Fórmula} & \textcolor{nerde}{\textbf{E.O.}} & \textcolor{nerde}{\textbf{Usual}} & \textcolor{nerde}{\textbf{IUPAC}} & \textcolor{nerde}{\textbf{Inglés}} \\
\midrule\endhead\endlastfoot
\ce{H^-} & 1$-$ & Hidruro & Hidruro(1$-$) & Hydride \\
\end{longtable}

\newpage

%---------- ANIONES COMPUESTOS ----------%
\phantomsection
\addcontentsline{toc}{section}{Aniones Compuestos}
\section*{\centering\fontsize{12}{18}\selectfont{\textcolor{blueunam}{\textbf{ANIONES COMPUESTOS}}}}

\vspace{0.3em}
\textbf{Cianuros y derivados ({\color{atomo}C, N, S, O})}

\begin{longtable}{C{2.2cm} L{4.8cm} >{\color{iupacgreen}}L{6.4cm} L{3.2cm}}
\rule{0pt}{2.8ex}\textbf{Fórmula} & \textcolor{nerde}{\textbf{Usual}} & \textcolor{nerde}{\textbf{IUPAC}} & \textcolor{nerde}{\textbf{Inglés}} \\
\midrule\endhead\endlastfoot
\ce{\ac{C}N^-} & Cianuro    & Cianuro(1$-$)    & Cyanide \\
\arrayrulecolor{separador}\midrule\arrayrulecolor{black}
\ce{\ac{S}CN^-} & Tiocianato & Tiocianato(1$-$) & Thiocyanate \\
\arrayrulecolor{separador}\midrule\arrayrulecolor{black}
\ce{O\ac{C}N^-} & Cianato    & Cianato(1$-$)    & Cyanate \\
\arrayrulecolor{separador}\midrule\arrayrulecolor{black}
\ce{\ac{C}NO^-} & Fulminato  & Fulminato(1$-$)  & Fulminate \\
\end{longtable}

\vspace{0.3em}
\textbf{Tiosales ({\color{atomo}S})}

\begin{longtable}{C{2.2cm} L{4.8cm} >{\color{iupacgreen}}L{6.4cm} L{3.2cm}}
\rule{0pt}{2.8ex}\textbf{Fórmula} & \textcolor{nerde}{\textbf{Usual}} & \textcolor{nerde}{\textbf{IUPAC}} & \textcolor{nerde}{\textbf{Inglés}} \\
\midrule\endhead\endlastfoot
\ce{\ac{S}2O3^{2-}} & Tiosulfato   & Tiosulfato(2$-$)      & Thiosulfate \\
\arrayrulecolor{separador}\midrule\arrayrulecolor{black}
\ce{\ac{S}2O8^{2-}} & Persulfato   & Peroxodisulfato(2$-$) & Persulfate \\
\arrayrulecolor{separador}\midrule\arrayrulecolor{black}
\ce{\ac{S}2O4^{2-}} & Ditionito    & Ditionito(2$-$)       & Dithionite \\
\arrayrulecolor{separador}\midrule\arrayrulecolor{black}
\ce{\ac{S}4O6^{2-}} & Tetrationato & Tetrationato(2$-$)    & Tetrathionate \\
\end{longtable}

\vspace{0.3em}
\textbf{Hidróxido y azidas}

\begin{longtable}{C{2.2cm} L{4.8cm} >{\color{iupacgreen}}L{6.4cm} L{3.2cm}}
\rule{0pt}{2.8ex}\textbf{Fórmula} & \textcolor{nerde}{\textbf{Usual}} & \textcolor{nerde}{\textbf{IUPAC}} & \textcolor{nerde}{\textbf{Inglés}} \\
\midrule\endhead\endlastfoot
\ce{\ac{O}H^-} & Hidróxido & Hidróxido(1$-$) & Hydroxide \\
\arrayrulecolor{separador}\midrule\arrayrulecolor{black}
\ce{\ac{N}3^-} & Azida     & Azida(1$-$)     & Azide \\
\arrayrulecolor{separador}\midrule\arrayrulecolor{black}
\ce{\ac{N}H2^-} & Amiduro   & Amiduro(1$-$)   & Amide \\
\ce{\ac{N}H^{2-}} & Imiduro   & Imiduro(2$-$)   & Imide \\
\end{longtable}

\newpage

%---------- SALES ÁCIDAS ----------%
\phantomsection
\addcontentsline{toc}{section}{Sales Ácidas}
\section*{\centering\fontsize{12}{18}\selectfont{\textcolor{blueunam}{\textbf{SALES ÁCIDAS}}}}

El prefijo \textbf{hidrógeno-} o \textbf{dihidrógeno-} en IUPAC indica cuántos H ionizables quedan sin sustituir.

\vspace{0.3em}
\textbf{Derivados del carbono ({\color{atomo}C})}

\begin{longtable}{C{2.2cm} >{\color{catred}}C{1.4cm} L{4.0cm} >{\color{iupacgreen}}L{6.2cm} L{2.8cm}}
\rule{0pt}{2.8ex}\textbf{Fórmula} & \textcolor{nerde}{\textbf{E.O.}} & \textcolor{nerde}{\textbf{Usual}} & \textcolor{nerde}{\textbf{IUPAC}} & \textcolor{nerde}{\textbf{Inglés}} \\
\midrule\endhead\endlastfoot
\ce{H\ac{C}O3^-} & 4+ & Bicarbonato & Hidrogenocarbonato(1$-$) & Hydrogen carbonate \\
\end{longtable}

\vspace{0.3em}
\textbf{Derivados del azufre ({\color{atomo}S})}

\begin{longtable}{C{2.2cm} >{\color{catred}}C{1.4cm} L{4.0cm} >{\color{iupacgreen}}L{6.2cm} L{2.8cm}}
\rule{0pt}{2.8ex}\textbf{Fórmula} & \textcolor{nerde}{\textbf{E.O.}} & \textcolor{nerde}{\textbf{Usual}} & \textcolor{nerde}{\textbf{IUPAC}} & \textcolor{nerde}{\textbf{Inglés}} \\
\midrule\endhead\endlastfoot
\ce{H\ac{S}O4^-} & 6+ & Bisulfato  & Hidrogenosulfato(1$-$)  & Hydrogen sulfate \\
\arrayrulecolor{separador}\midrule\arrayrulecolor{black}
\ce{H\ac{S}O3^-} & 4+ & Bisulfito  & Hidrogenosulfito(1$-$)  & Hydrogen sulfite \\
\end{longtable}

\vspace{0.3em}
\textbf{Derivados del fósforo ({\color{atomo}P})}

\begin{longtable}{C{2.2cm} >{\color{catred}}C{1.4cm} L{4.0cm} >{\color{iupacgreen}}L{6.2cm} L{2.8cm}}
\rule{0pt}{2.8ex}\textbf{Fórmula} & \textcolor{nerde}{\textbf{E.O.}} & \textcolor{nerde}{\textbf{Usual}} & \textcolor{nerde}{\textbf{IUPAC}} & \textcolor{nerde}{\textbf{Inglés}} \\
\midrule\endhead\endlastfoot
\ce{H2\ac{P}O4^-} & 5+ & Fosfato diácido   & Dihidrogenofosfato(1$-$)  & Dihydrogen phosphate \\
\ce{H\ac{P}O4^{2-}} & 5+ & Fosfato monoácido & Hidrogenofosfato(2$-$)    & Hydrogen phosphate \\
\arrayrulecolor{separador}\midrule\arrayrulecolor{black}
\ce{H2\ac{P}O3^-} & 3+ & Fosfito diácido   & Dihidrogenofosfito(1$-$)  & Dihydrogen phosphite \\
\ce{H\ac{P}O3^{2-}} & 3+ & Fosfito monoácido & Hidrogenofosfito(2$-$)    & Hydrogen phosphite \\
\end{longtable}

\newpage

%---------- ÁCIDOS OXÁCIDOS ----------%
\phantomsection
\addcontentsline{toc}{section}{Ácidos Oxácidos}
\section*{\centering\fontsize{12}{18}\selectfont{\textcolor{blueunam}{\textbf{ÁCIDOS OXÁCIDOS}}}}

Usual: \textbf{-ato} $\to$ \textbf{-ico}; \quad \textbf{-ito} $\to$ \textbf{-oso}.\quad
IUPAC: ácido + $n\times$oxo + raíz + -ico + (E.O. en romano).

\vspace{0.3em}
\textbf{Oxácidos: Halógenos ({\color{atomo}Cl, Br, I})}

\begin{longtable}{C{2.2cm} >{\color{catred}}C{1.4cm} L{4.0cm} >{\color{iupacgreen}}L{6.2cm} L{2.8cm}}
\rule{0pt}{2.8ex}\textbf{Fórmula} & \textcolor{nerde}{\textbf{E.O.}} & \textcolor{nerde}{\textbf{Usual}} & \textcolor{nerde}{\textbf{IUPAC}} & \textcolor{nerde}{\textbf{Inglés}} \\
\midrule\endhead\endlastfoot
\ce{H\ac{Cl}O} & 1+ & Ácido hipocloroso  & Ácido oxoclórico(I)        & Hypochlorous acid \\
\ce{H\ac{Cl}O2} & 3+ & Ácido cloroso      & Ácido dioxoclórico(III)    & Chlorous acid \\
\ce{H\ac{Cl}O3} & 5+ & Ácido clórico      & Ácido trioxoclórico(V)     & Chloric acid \\
\ce{H\ac{Cl}O4} & 7+ & Ácido perclórico   & Ácido tetraoxoclórico(VII) & Perchloric acid \\
\arrayrulecolor{separador}\midrule\arrayrulecolor{black}
\ce{H\ac{Br}O} & 1+ & Ácido hipobromoso  & Ácido oxobrómico(I)        & Hypobromous acid \\
\ce{H\ac{Br}O2} & 3+ & Ácido bromoso      & Ácido dioxobrómico(III)    & Bromous acid \\
\ce{H\ac{Br}O3} & 5+ & Ácido brómico      & Ácido trioxobrómico(V)     & Bromic acid \\
\ce{H\ac{Br}O4} & 7+ & Ácido perbrómico   & Ácido tetraoxobrómico(VII) & Perbromic acid \\
\arrayrulecolor{separador}\midrule\arrayrulecolor{black}
\ce{H\ac{I}O} & 1+ & Ácido hipoyodoso   & Ácido oxoyódico(I)         & Hypoiodous acid \\
\ce{H\ac{I}O2} & 3+ & Ácido yodoso       & Ácido dioxoyódico(III)     & Iodous acid \\
\ce{H\ac{I}O3} & 5+ & Ácido yódico       & Ácido trioxoyódico(V)      & Iodic acid \\
\ce{H\ac{I}O4} & 7+ & Ácido peryódico    & Ácido tetraoxoyódico(VII)  & Periodic acid \\
\end{longtable}

\vspace{0.3em}
\textbf{Oxácidos: Anfígenos ({\color{atomo}S, Se, Te})}

\begin{longtable}{C{2.2cm} >{\color{catred}}C{1.4cm} L{4.0cm} >{\color{iupacgreen}}L{6.2cm} L{2.8cm}}
\rule{0pt}{2.8ex}\textbf{Fórmula} & \textcolor{nerde}{\textbf{E.O.}} & \textcolor{nerde}{\textbf{Usual}} & \textcolor{nerde}{\textbf{IUPAC}} & \textcolor{nerde}{\textbf{Inglés}} \\
\midrule\endhead\endlastfoot
\ce{H2\ac{S}O3} & 4+ & Ácido sulfuroso   & Ácido trioxosulfúrico(IV)   & Sulfurous acid \\
\ce{H2\ac{S}O4} & 6+ & Ácido sulfúrico   & Ácido tetraoxosulfúrico(VI) & Sulfuric acid \\
\arrayrulecolor{separador}\midrule\arrayrulecolor{black}
\ce{H2\ac{Se}O3} & 4+ & Ácido selenoso    & Ácido trioxoselénico(IV)    & Selenous acid \\
\ce{H2\ac{Se}O4} & 6+ & Ácido selénico    & Ácido tetraoxoselénico(VI)  & Selenic acid \\
\arrayrulecolor{separador}\midrule\arrayrulecolor{black}
\ce{H2\ac{Te}O3} & 4+ & Ácido teluroso    & Ácido trioxotelúrico(IV)    & Tellurous acid \\
\ce{H2\ac{Te}O4} & 6+ & Ácido telúrico    & Ácido tetraoxotelúrico(VI)  & Telluric acid \\
\end{longtable}

\vspace{0.3em}
\textbf{Oxácidos: Nitrogenoides ({\color{atomo}N, P, As})}

\begin{longtable}{C{2.2cm} >{\color{catred}}C{1.4cm} L{4.0cm} >{\color{iupacgreen}}L{6.2cm} L{2.8cm}}
\rule{0pt}{2.8ex}\textbf{Fórmula} & \textcolor{nerde}{\textbf{E.O.}} & \textcolor{nerde}{\textbf{Usual}} & \textcolor{nerde}{\textbf{IUPAC}} & \textcolor{nerde}{\textbf{Inglés}} \\
\midrule\endhead\endlastfoot
\ce{H\ac{N}O2} & 3+ & Ácido nitroso      & Ácido dioxonítrico(III)      & Nitrous acid \\
\ce{H\ac{N}O3} & 5+ & Ácido nítrico      & Ácido trioxonítrico(V)       & Nitric acid \\
\arrayrulecolor{separador}\midrule\arrayrulecolor{black}
\ce{H3\ac{P}O3} & 3+ & Ácido fosforoso    & Ácido trioxofosfórico(III)   & Phosphorous acid \\
\ce{H3\ac{P}O4} & 5+ & Ácido fosfórico    & Ácido tetraoxofosfórico(V)   & Phosphoric acid \\
\arrayrulecolor{separador}\midrule\arrayrulecolor{black}
\ce{H3\ac{As}O3} & 3+ & Ácido arsenioso    & Ácido trioxoarsénico(III)    & Arsenious acid \\
\ce{H3\ac{As}O4} & 5+ & Ácido arsénico     & Ácido tetraoxoarsénico(V)    & Arsenic acid \\
\end{longtable}

\vspace{0.3em}
\textbf{Carbonoideos ({\color{atomo}C, Si})}

\begin{longtable}{C{2.2cm} >{\color{catred}}C{1.4cm} L{4.0cm} >{\color{iupacgreen}}L{6.2cm} L{2.8cm}}
\rule{0pt}{2.8ex}\textbf{Fórmula} & \textcolor{nerde}{\textbf{E.O.}} & \textcolor{nerde}{\textbf{Usual}} & \textcolor{nerde}{\textbf{IUPAC}} & \textcolor{nerde}{\textbf{Inglés}} \\
\midrule\endhead\endlastfoot
\ce{H2\ac{C}O3} & 4+ & Ácido carbónico & Ácido trioxocarbónico(IV)  & Carbonic acid \\
\arrayrulecolor{separador}\midrule\arrayrulecolor{black}
\ce{H2\ac{Si}O3} & 4+ & Ácido silícico  & Ácido trioxosilícico(IV)   & Silicic acid \\
\end{longtable}

\vspace{0.3em}
\textbf{Oxácidos: Metales de transición ({\color{atomo}Cr, Mn, V, Mo, W})}

\begin{longtable}{C{2.2cm} >{\color{catred}}C{1.4cm} L{4.0cm} >{\color{iupacgreen}}L{6.2cm} L{2.8cm}}
\rule{0pt}{2.8ex}\textbf{Fórmula} & \textcolor{nerde}{\textbf{E.O.}} & \textcolor{nerde}{\textbf{Usual}} & \textcolor{nerde}{\textbf{IUPAC}} & \textcolor{nerde}{\textbf{Inglés}} \\
\midrule\endhead\endlastfoot
\ce{H2\ac{Cr}O4} & 6+ & Ácido crómico      & Ácido tetraoxocrómico(VI)     & Chromic acid \\
\ce{H2\ac{Cr}2O7} & 6+ & Ácido dicrómico    & Ácido heptaoxodicrómico(VI)   & Dichromic acid \\
\arrayrulecolor{separador}\midrule\arrayrulecolor{black}
\ce{H\ac{Mn}O4} & 7+ & Ácido permangánico & Ácido tetraoxomangánico(VII)  & Permanganic acid \\
\ce{H2\ac{Mn}O4} & 6+ & Ácido mangánico    & Ácido tetraoxomangánico(VI)   & Manganic acid \\
\arrayrulecolor{separador}\midrule\arrayrulecolor{black}
\ce{H\ac{V}O3} & 5+ & Ácido vanádico     & Ácido trioxovanádico(V)       & Vanadic acid \\
\arrayrulecolor{separador}\midrule\arrayrulecolor{black}
\ce{H2\ac{Mo}O4} & 6+ & Ácido molíbdico    & Ácido tetraoxomolíbdico(VI)   & Molybdic acid \\
\arrayrulecolor{separador}\midrule\arrayrulecolor{black}
\ce{H2\ac{W}O4} & 6+ & Ácido túngstico    & Ácido tetraoxotúngstico(VI)   & Tungstic acid \\
\end{longtable}

\newpage

%---------- ÁCIDOS HIDRÁCIDOS ----------%
\phantomsection
\addcontentsline{toc}{section}{Ácidos Hidrácidos}
\section*{\centering\fontsize{12}{18}\selectfont{\textcolor{blueunam}{\textbf{ÁCIDOS HIDRÁCIDOS}}}}

Compuesto puro (gas): anión + ``de hidrógeno'' (IUPAC). En solución acuosa: ácido + raíz + \textbf{-hídrico} (Usual). El \ce{H+} real en solución es \ce{H3O+}; se escribe \ce{H+} como simplificación.

\vspace{0.3em}
\textbf{Halógenos ({\color{atomo}F, Cl, Br, I})}

\begin{longtable}{C{1.6cm} C{1.8cm} >{\color{catred}}C{1.2cm} L{3.6cm} >{\color{iupacgreen}}L{5.0cm} L{3.4cm}}
\rule{0pt}{2.8ex}\textbf{F.(g)} & \textbf{F.(ac)} & \textcolor{nerde}{\textbf{E.O.}} & \textcolor{nerde}{\textbf{Usual (ac)}} & \textcolor{nerde}{\textbf{IUPAC (puro)}} & \textcolor{nerde}{\textbf{Inglés}} \\
\midrule\endhead\endlastfoot
\ce{H\ac{F}} & \ce{H\ac{F}(aq)}  & 1$-$ & Ácido fluorhídrico  & Fluoruro de hidrógeno  & Hydrofluoric acid \\
\arrayrulecolor{separador}\midrule\arrayrulecolor{black}
\ce{H\ac{Cl}} & \ce{H\ac{Cl}(aq)} & 1$-$ & Ácido clorhídrico   & Cloruro de hidrógeno   & Hydrochloric acid \\
\arrayrulecolor{separador}\midrule\arrayrulecolor{black}
\ce{H\ac{Br}} & \ce{H\ac{Br}(aq)} & 1$-$ & Ácido bromhídrico   & Bromuro de hidrógeno   & Hydrobromic acid \\
\arrayrulecolor{separador}\midrule\arrayrulecolor{black}
\ce{H\ac{I}} & \ce{H\ac{I}(aq)}  & 1$-$ & Ácido yodhídrico    & Yoduro de hidrógeno    & Hydroiodic acid \\
\end{longtable}

\vspace{0.3em}
\textbf{Anfígenos ({\color{atomo}S, Se, Te})}

\begin{longtable}{C{1.6cm} C{1.8cm} >{\color{catred}}C{1.2cm} L{3.6cm} >{\color{iupacgreen}}L{5.0cm} L{3.4cm}}
\rule{0pt}{2.8ex}\textbf{F.(g)} & \textbf{F.(ac)} & \textcolor{nerde}{\textbf{E.O.}} & \textcolor{nerde}{\textbf{Usual (ac)}} & \textcolor{nerde}{\textbf{IUPAC (puro)}} & \textcolor{nerde}{\textbf{Inglés}} \\
\midrule\endhead\endlastfoot
\ce{H2\ac{S}}  & \ce{H2\ac{S}(aq)}  & 2$-$ & Ácido sulfhídrico   & Sulfuro de hidrógeno   & Hydrosulfuric acid \\
\arrayrulecolor{separador}\midrule\arrayrulecolor{black}
\ce{H2\ac{Se}} & \ce{H2\ac{Se}(aq)} & 2$-$ & Ácido selenhídrico  & Seleniuro de hidrógeno & Hydroselenic acid \\
\arrayrulecolor{separador}\midrule\arrayrulecolor{black}
\ce{H2\ac{Te}} & \ce{H2\ac{Te}(aq)} & 2$-$ & Ácido telurhídrico  & Telururo de hidrógeno  & Hydrotelluric acid \\
\end{longtable}

\vspace{0.3em}
\textbf{Nitrogenoides y carbonoideos ({\color{atomo}N, C})}

\begin{longtable}{C{1.6cm} C{1.8cm} >{\color{catred}}C{1.2cm} L{3.6cm} >{\color{iupacgreen}}L{5.0cm} L{3.4cm}}
\rule{0pt}{2.8ex}\textbf{F.(g)} & \textbf{F.(ac)} & \textcolor{nerde}{\textbf{E.O.}} & \textcolor{nerde}{\textbf{Usual (ac)}} & \textcolor{nerde}{\textbf{IUPAC (puro)}} & \textcolor{nerde}{\textbf{Inglés}} \\
\midrule\endhead\endlastfoot
\ce{H\ac{N}3} & \ce{H\ac{N}3(aq)} & 1$-$ & Ácido hidrazoico  & Azida de hidrógeno   & Hydrazoic acid \\
\arrayrulecolor{separador}\midrule\arrayrulecolor{black}
\ce{H\ac{C}N}  & \ce{H\ac{C}N(aq)}  & 1$-$ & Ácido cianhídrico & Cianuro de hidrógeno & Hydrocyanic acid \\
\end{longtable}

\newpage

%---------- CATIONES ----------%
\begin{landscape}
\phantomsection
\addcontentsline{toc}{section}{Cationes Comunes}
\section*{\centering\fontsize{12}{18}\selectfont{\textcolor{blueunam}{\textbf{CATIONES COMUNES}}}}

\vspace*{0.2cm}
\begin{center}
{\large\textbf{\textcolor{blueunam}{Cationes}}}
\end{center}

\vspace{0.3em}
\setlength{\tabcolsep}{3pt}
\renewcommand{\arraystretch}{1.1}
\small

% --- INICIO DE LA TABLA DE CATIONES ---
% 1. Forzamos el interlineado a 1.0 para que las filas no sean tan altas
\renewcommand{\arraystretch}{1.0}

% 2. Usamos \parbox con altura fija (1.05cm). Esto evita que el color se desborde y tape las líneas.
\renewcommand{\catcell}[4]{%
    \cellcolor{#1}%
    \parbox[c][1.05cm][c]{1.25cm}{%
        \centering
        {\scriptsize\textbf{#2}}\par\vspace{0.5pt}%
        {\tiny\color{gray}#3}\par\vspace{0.5pt}%
        {\tiny\color{blueunam}#4}%
    }%
}

% 3. Aseguramos líneas negras y del grosor correcto
\arrayrulecolor{black}
\setlength{\arrayrulewidth}{0.4pt}

\noindent
\begin{tabular}{|c|C{1.35cm}|C{1.35cm}|C{1.35cm}|C{1.35cm}|C{1.35cm}|C{1.35cm}|C{1.35cm}|C{1.35cm}|C{1.35cm}|C{1.35cm}|C{1.35cm}|C{1.35cm}|C{1.35cm}|C{1.35cm}|C{1.35cm}|}
\hline
\rowcolor{blueunam}
\textcolor{white}{\scriptsize\textbf{Per.}} &
\textcolor{white}{\scriptsize\textbf{G1}} &
\textcolor{white}{\scriptsize\textbf{G2}} &
\textcolor{white}{\scriptsize\textbf{G3}} &
\textcolor{white}{\scriptsize\textbf{G4}} &
\textcolor{white}{\scriptsize\textbf{G5}} &
\textcolor{white}{\scriptsize\textbf{G6}} &
\textcolor{white}{\scriptsize\textbf{G7}} &
\textcolor{white}{\scriptsize\textbf{G8}} &
\textcolor{white}{\scriptsize\textbf{G9}} &
\textcolor{white}{\scriptsize\textbf{G10}} &
\textcolor{white}{\scriptsize\textbf{G11}} &
\textcolor{white}{\scriptsize\textbf{G12}} &
\textcolor{white}{\scriptsize\textbf{G13}} &
\textcolor{white}{\scriptsize\textbf{G14}} &
\textcolor{white}{\scriptsize\textbf{G15}} \\
\hline
{\scriptsize\textbf{1}} &
\catcell{col1}{\ce{H+}}{}{1+} &
& & & & & & & & & & & & & \\
\hline
{\scriptsize\textbf{2}} &
\catcell{col1}{\ce{Li+}}{}{1+} &
\catcell{col2}{\ce{Be^{2+}}}{}{2+} &
& & & & & & & & & & 
\catcell{pos}{\ce{B^{3+}}}{}{3+} &
& \\
\hline
{\scriptsize\textbf{3}} &
\catcell{col1}{\ce{Na+}}{}{1+} &
\catcell{col2}{\ce{Mg^{2+}}}{}{2+} &
& & & & & & & & & & 
\catcell{pos}{\ce{Al^{3+}}}{}{3+} &
& \\
\hline
{\scriptsize\textbf{4}} &
\catcell{col1}{\ce{K+}}{}{1+} &
\catcell{col2}{\ce{Ca^{2+}}}{}{2+} &
\catcell{trans}{\ce{Sc^{3+}}}{}{3+} &
\catcell{trans}{\ce{Ti^{4+}}}{}{4+} &
\catcell{trans}{\ce{V^{5+}}}{}{5+} &
\catcell{trans}{\ce{Cr^{3+}}}{}{3+} &
\catcell{trans}{\ce{Mn^{2+}}}{}{2+} &
\catcell{trans}{\ce{Fe^{2+}/Fe^{3+}}}{}{2+,3+} &
\catcell{trans}{\ce{Co^{2+}}}{}{2+} &
\catcell{trans}{\ce{Ni^{2+}}}{}{2+} &
\catcell{trans}{\ce{Cu+/Cu^{2+}}}{}{1+,2+} &
\catcell{trans}{\ce{Zn^{2+}}}{}{2+} &
\catcell{pos}{\ce{Ga^{3+}}}{}{3+} &
\catcell{pos}{\ce{Ge^{4+}}}{}{4+} &
\catcell{pos}{\ce{As^{3+}}}{}{3+} \\
\hline
{\scriptsize\textbf{5}} &
\catcell{col1}{\ce{Rb+}}{}{1+} &
\catcell{col2}{\ce{Sr^{2+}}}{}{2+} &
\catcell{trans}{\ce{Y^{3+}}}{}{3+} &
\catcell{trans}{\ce{Zr^{4+}}}{}{4+} &
\catcell{trans}{\ce{Nb^{5+}}}{}{5+} &
\catcell{trans2}{\ce{Mo^{6+}}}{}{6+} &
& & &
\catcell{trans2}{\ce{Pd^{2+}}}{}{2+} &
\catcell{trans2}{\ce{Ag+}}{}{1+} &
\catcell{trans2}{\ce{Cd^{2+}}}{}{2+} &
\catcell{pos}{\ce{In^{3+}}}{}{3+} &
\catcell{pos}{\ce{Sn^{2+}/Sn^{4+}}}{}{2+,4+} &
\catcell{pos}{\ce{Sb^{3+}}}{}{3+} \\
\hline
{\scriptsize\textbf{6}} &
\catcell{col1}{\ce{Cs+}}{}{1+} &
\catcell{col2}{\ce{Ba^{2+}}}{}{2+} &
\catcell{trans}{\ce{La^{3+}}}{}{3+} &
\catcell{trans}{\ce{Hf^{4+}}}{}{4+} &
\catcell{trans}{\ce{Ta^{5+}}}{}{5+} &
\catcell{trans2}{\ce{W^{6+}}}{}{6+} &
& & &
\catcell{trans2}{\ce{Pt^{2+}/Pt^{4+}}}{}{2+,4+} &
\catcell{trans2}{\ce{Au+/Au^{3+}}}{}{1+,3+} &
\catcell{trans2}{\scalebox{0.85}{\ce{Hg2^{2+}/Hg^{2+}}}}{}{(I)*,(II)} &
\catcell{pos}{\scalebox{0.9}{\ce{Tl+/Tl^{3+}}}}{}{1+,3+} &
\catcell{pos}{\ce{Pb^{2+}/Pb^{4+}}}{}{2+,4+} &
\catcell{pos}{\ce{Bi^{3+}}}{}{3+} \\
\hline
{\scriptsize\textbf{7}} &
\catcell{col1}{\ce{Fr+}}{}{1+} &
\catcell{col2}{\ce{Ra^{2+}}}{}{2+} &
& & & & & & & & & & & & \\
\hline
\end{tabular}

% Restauramos el interlineado original para que no afecte al resto del documento
\renewcommand{\arraystretch}{1.3}

\vspace{0.4em}
{\small
*\ce{Hg2^{2+}}: catión impar, no se simplifica a \ce{Hg+}. Los cationes pares sí se simplifican.\\
Cr también forma Cr$^{2+}$; Fe, Co, Ni también tienen estados 3+; detalle en tabla siguiente.
}
% --- FIN DE LA TABLA DE CATIONES ---

\vspace{0.8em}

\textbf{Cationes fuera de la tabla periódica}

\begin{longtable}{C{2.2cm} >{\color{catred}}C{1.4cm} L{4.0cm} >{\color{iupacgreen}}L{6.2cm} L{2.8cm}}
\rule{0pt}{2.8ex}\textbf{Fórmula} & \textcolor{nerde}{\textbf{E.O.}} & \textcolor{nerde}{\textbf{Usual}} & \textcolor{nerde}{\textbf{IUPAC}} & \textcolor{nerde}{\textbf{Inglés}} \\
\midrule\endhead\endlastfoot
\ce{\ac{N}H4^+} & 1+ & Amonio   & Amonio(1+)   & Ammonium(1+) \\
\arrayrulecolor{separador}\midrule\arrayrulecolor{black}
\ce{H3\ac{O}^+} & 1+ & Hidronio & Hidronio(1+) & Hydronium(1+) \\
\end{longtable}

\end{landscape}

\renewcommand{\arraystretch}{1.3}
\setlength{\tabcolsep}{5pt}

\newpage

\restoregeometry
%---------- REGLAS DE NOMENCLATURA ----------%
\phantomsection
\addcontentsline{toc}{section}{Reglas de Nomenclatura: Formación de Fórmulas y Sales}
\section*{\centering\fontsize{12}{18}\selectfont{\textcolor{blueunam}{\textbf{REGLAS DE NOMENCLATURA: FORMACIÓN DE FÓRMULAS Y SALES}}}}

\subsection*{\textcolor{blueunam}{Principio de neutralidad y formación de fórmulas}}

La carga total del compuesto debe ser \textbf{cero}. Procedimiento: (1) escribir las cargas, (2) intercambiarlas como subíndices sin signos, (3) simplificar si ambos son divisibles, (4) anión poliatómico repetido va entre paréntesis.

\begin{longtable}{C{2.4cm} C{2.4cm} C{2.8cm} C{2.4cm}}
\toprule
\textbf{Catión} & \textbf{Anión} & \textbf{Intercambio} & \textbf{Fórmula} \\
\midrule\endfirsthead\toprule
\textbf{Catión} & \textbf{Anión} & \textbf{Intercambio} & \textbf{Fórmula} \\
\midrule\endhead\bottomrule\endlastfoot
\ce{Na+} (1+)     & \ce{Cl-} (1$-$)       & 1 y 1 & \ce{NaCl} \\
\ce{Ca^{2+}} (2+) & \ce{Cl-} (1$-$)       & 2 y 1 & \ce{CaCl2} \\
\ce{Al^{3+}} (3+) & \ce{O^{2-}} (2$-$)    & 3 y 2 & \ce{Al2O3} \\
\ce{Fe^{3+}} (3+) & \ce{SO4^{2-}} (2$-$)  & 3 y 2 & \ce{Fe2(SO4)3} \\
\ce{NH4+} (1+)    & \ce{NO3-} (1$-$)      & 1 y 1 & \ce{NH4NO3} \\
\ce{Ca^{2+}} (2+) & \ce{OH-} (1$-$)       & 2 y 1 & \ce{Ca(OH)2} \\
\bottomrule
\end{longtable}

\subsection*{\textcolor{blueunam}{Nomenclatura de sales}}

Nombre Usual: \textbf{anión + ``de'' + catión (con romano si aplica)}.

\begin{longtable}{C{2.0cm} L{2.6cm} L{2.6cm} L{4.4cm}}
\toprule
\textbf{Tipo} & \textbf{Anión} & \textbf{Catión} & \textbf{Nombre Usual} \\
\midrule\endfirsthead\toprule
\textbf{Tipo} & \textbf{Anión} & \textbf{Catión} & \textbf{Nombre Usual} \\
\midrule\endhead\bottomrule\endlastfoot
Binaria     & \ce{Cl-} (cloruro)       & \ce{Na+} (sodio)          & Cloruro de sodio \\
Binaria     & \ce{O^{2-}} (óxido)      & \ce{Fe^{2+}} (hierro(II)) & Óxido de hierro(II) \\
Ternaria    & \ce{SO4^{2-}} (sulfato)  & \ce{Cu^{2+}} (cobre(II))  & Sulfato de cobre(II) \\
Ternaria    & \ce{NO3-} (nitrato)      & \ce{Ag+} (plata)          & Nitrato de plata \\
Ácida       & \ce{HCO3-} (bicarbonato) & \ce{Na+} (sodio)          & Bicarbonato de sodio \\
Hidróxido   & \ce{OH-} (hidróxido)     & \ce{Ca^{2+}} (calcio)     & Hidróxido de calcio \\
\bottomrule
\end{longtable}

\subsection*{\textcolor{blueunam}{Hidróxidos}}

\begin{longtable}{C{2.8cm} L{4.2cm}}
\toprule
\textbf{Fórmula} & \textbf{Nombre Usual} \\
\midrule\endfirsthead\toprule
\textbf{Fórmula} & \textbf{Nombre Usual} \\
\midrule\endhead\bottomrule\endlastfoot
\ce{NaOH}          & Hidróxido de sodio \\
\ce{Ca(OH)2}  & Hidróxido de calcio \\
\ce{Fe(OH)3}  & Hidróxido de hierro(III) \\
\ce{Al(OH)3}  & Hidróxido de aluminio \\
\bottomrule
\end{longtable}

\subsection*{\textcolor{blueunam}{Óxidos}}

\begin{longtable}{C{2.0cm} L{2.8cm} L{4.0cm} L{3.6cm}}
\toprule
\textbf{Tipo} & \textbf{Fórmula} & \textbf{Nombre Usual} & \textbf{IUPAC} \\
\midrule\endfirsthead\toprule
\textbf{Tipo} & \textbf{Fórmula} & \textbf{Nombre Usual} & \textbf{IUPAC} \\
\midrule\endhead\bottomrule\endlastfoot
Básico & \ce{FeO}        & Óxido de hierro(II)  & Óxido de hierro(2+) \\
Básico & \ce{Fe2O3} & Óxido de hierro(III) & Óxido de hierro(3+) \\
Ácido  & \ce{CO2}   & Óxido de carbono(IV) & Óxido de carbono(4+) \\
Ácido  & \ce{SO3}   & Óxido de azufre(VI)  & Óxido de azufre(6+) \\
\bottomrule
\end{longtable}

\subsection*{\textcolor{blueunam}{Peróxidos}}

\begin{longtable}{C{2.4cm} L{3.6cm} L{3.6cm}}
\toprule
\textbf{Fórmula} & \textbf{Nombre Usual} & \textbf{IUPAC} \\
\midrule\endfirsthead\toprule
\textbf{Fórmula} & \textbf{Nombre Usual} & \textbf{IUPAC} \\
\midrule\endhead\bottomrule\endlastfoot
\ce{H2O2}  & Peróxido de hidrógeno & Peróxido de hidrógeno \\
\ce{Na2O2} & Peróxido de sodio     & Peróxido de sodio \\
\ce{BaO2}  & Peróxido de bario     & Peróxido de bario \\
\bottomrule
\end{longtable}

El ion es \ce{O2^{2-}}, no \ce{O^{2-}}. El oxígeno está en estado $-1$, no $-2$.

\subsection*{\textcolor{blueunam}{Hidruros}}

\begin{longtable}{C{2.0cm} L{2.8cm} L{5.4cm}}
\toprule
\textbf{Tipo} & \textbf{Fórmula} & \textbf{Nombre} \\
\midrule\endfirsthead\toprule
\textbf{Tipo} & \textbf{Fórmula} & \textbf{Nombre} \\
\midrule\endhead\bottomrule\endlastfoot
Metálico    & \ce{NaH}       & Hidruro de sodio (Usual e IUPAC) \\
Metálico    & \ce{CaH2} & Hidruro de calcio \\
No metálico & \ce{NH3}  & Amoníaco (nombre especial) \\
No metálico & \ce{H2O}  & Agua (nombre especial) \\
\bottomrule
\end{longtable}

\subsection*{\textcolor{blueunam}{Resumen de los tres sistemas}}

\begin{longtable}{L{2.4cm} L{4.6cm} L{4.8cm}}
\toprule
\textbf{Sistema} & \textbf{Cómo nombra} & \textbf{Ejemplo: \ce{Fe2(SO4)3}} \\
\midrule\endfirsthead\toprule
\textbf{Sistema} & \textbf{Cómo nombra} & \textbf{Ejemplo: \ce{Fe2(SO4)3}} \\
\midrule\endhead\bottomrule\endlastfoot
Tradicional & Prefijos y sufijos para catión y anión      & Sulfato férrico \\
Usual       & Tradicional para anión + romano para catión  & Sulfato de hierro(III) \\
IUPAC       & Carga entre paréntesis con arábigos          & Sulfato(2$-$) de hierro(3+) \\
\bottomrule
\end{longtable}

\textbf{Regla de oro:} anión primero, catión después, unidos por ``de''. El anión lleva sufijo -uro, -ato, -ito u -óxido según su tipo. El catión lleva número romano solo si tiene más de un estado de oxidación (en Usual). En IUPAC, aniones y cationes llevan su carga entre paréntesis; en oxácidos el E.O. va en romano.

%---------- Final del documento ----------%
\end{document}